\chapter{量子光学}
\label{ch:quantum-optics}

半经典原子--光耦合把电磁场当作给定函数，因此能描述受驱吸收和 Rabi 振荡；可是一旦外场关闭，这个哈密顿量不能让激发原子向真空自发放出光子。缺少的自由度正是电磁场自身的量子涨落。

腔中的每个 Maxwell 正常模都具有谐振子的能量形式，因此可以逐模量子化。升降算符定义后立即进入电偶极矩阵元；对连续末态求和便得到自发辐射率，Green 并矢再把自由空间结果推广到非均匀环境。Fock、相干和压缩光场的统计性质留到下一章。
\section{自伴 Maxwell 正常模与场的 Hilbert 空间}

量子化的首要对象不是“平面波光子”，而是经典 Maxwell 方程的正交正常模。先考虑无耗散、非磁性但允许空间非均匀的介质，$\epsilon_r(\mathbf r)>0$，$\mu=\mu_0$。在无自由电荷的广义 Coulomb 条件下，横向正常模满足
\begin{equation}
\nabla\times\nabla\times\mathbf f_\lambda(\mathbf r)
=\frac{\omega_\lambda^2}{c^2}\epsilon_r(\mathbf r)\mathbf f_\lambda(\mathbf r),
\qquad
\nabla\cdot[\epsilon_r\mathbf f_\lambda]=0.
\label{eq:maxwell-normal-mode-eigenproblem}
\end{equation}
自然内积是
\begin{equation}
\langle\mathbf f_\lambda,\mathbf f_{\lambda'}\rangle_\epsilon
=\int\dd^3r\,\epsilon_r(\mathbf r)
\mathbf f_\lambda^*(\mathbf r)\cdot\mathbf f_{\lambda'}(\mathbf r)
=\delta_{\lambda\lambda'}.
\label{eq:maxwell-weighted-inner-product}
\end{equation}
在使边界通量项消失的边界条件下，$\epsilon_r^{-1}\nabla\times\nabla\times$ 对该加权内积自伴，因此不同 $\omega_\lambda^2$ 的模式正交，场能可以分解为独立谐振子。真空平面波、腔模、光子晶体 Bloch 模和波导离散横模都只是这一本征问题的不同边界条件或材料特例。

\begin{proof}[解]
取两个满足相同边界条件的横向场 \(\mathbf f,\mathbf g\)。由矢量恒等式
\begin{equation*}
 \nabla\cdot(\mathbf A\times\mathbf B)
 =\mathbf B\cdot(\nabla\times\mathbf A)
 -\mathbf A\cdot(\nabla\times\mathbf B),
\end{equation*}
取 \(\mathbf A=\mathbf f^*\)、\(\mathbf B=\nabla\times\mathbf g\) 与 \(\mathbf A=(\nabla\times\mathbf f)^*\)、\(\mathbf B=\mathbf g\) 分别展开并相减，
\begin{align*}
 \mathbf f^*\cdot(\nabla\times\nabla\times\mathbf g)
 &-(\nabla\times\nabla\times\mathbf f)^*\cdot\mathbf g\\
 &=\nabla\cdot\left[\mathbf f^*\times(\nabla\times\nabla\times\mathbf g)\right]
 -\nabla\cdot\left[(\nabla\times\mathbf f)^*\times(\nabla\times\mathbf g)\right]\\
 &\quad+\nabla\cdot\left[(\nabla\times\nabla\times\mathbf f)^*\times\mathbf g\right]
 -\nabla\cdot\left[(\nabla\times\mathbf f)^*\times(\nabla\times\mathbf g)\right]\\
 &=\nabla\cdot\mathbf J_{fg},
\end{align*}
对体积积分并由散度定理，
\begin{equation*}
 \int\dd^3r\,\mathbf f^*\cdot(\nabla\times\nabla\times\mathbf g)
 -\int\dd^3r\,(\nabla\times\nabla\times\mathbf f)^*\cdot\mathbf g
 =\oint\dd\mathbf S\cdot\mathbf J_{fg},
\end{equation*}
其中 \(\mathbf J_{fg}\) 是由 \(\mathbf f,\mathbf g\) 及其旋度构成的边界通量。周期边界使切向分量匹配、理想导体边界使切向 \(\mathbf E\) 为零、无穷远足够快衰减使面积分趋于零，三者都使右端为零。

令 \(\mathcal M=\epsilon_r^{-1}\nabla\times\nabla\times\)，加权内积\eqnrefs{eq:maxwell-weighted-inner-product}给出
\begin{align*}
 \langle\mathbf f,\mathcal M\mathbf g\rangle_\epsilon
 &=\int\dd^3r\,\epsilon_r\mathbf f^*\cdot
 \left(\epsilon_r^{-1}\nabla\times\nabla\times\mathbf g\right)\\
 &=\int\dd^3r\,\mathbf f^*\cdot(\nabla\times\nabla\times\mathbf g)\\
 &=\int\dd^3r\,(\nabla\times\nabla\times\mathbf f)^*\cdot\mathbf g\\
 &=\int\dd^3r\,\epsilon_r
 \left(\epsilon_r^{-1}\nabla\times\nabla\times\mathbf f\right)^*\cdot\mathbf g\\
 &=\langle\mathcal M\mathbf f,\mathbf g\rangle_\epsilon,
\end{align*}
其中第一步用了 \(\epsilon_r\cdot\epsilon_r^{-1}=1\)，第三步用了边界通量为零的分部积分。因此 \(\mathcal M\) 在加权内积下自伴。

由自伴性，\eqnrefs{eq:maxwell-normal-mode-eigenproblem}的本征值 \(\omega_\lambda^2/c^2\) 为实谱：对 \(\mathcal M\mathbf f_\lambda=\frac{\omega_\lambda^2}{c^2}\mathbf f_\lambda\) 取复共轭并利用 \(\mathcal M^\dagger=\mathcal M\) 即得 \((\omega_\lambda^2)^*=\omega_\lambda^2\)。对不同本征值 \(\omega_\lambda^2\ne\omega_{\lambda'}^2\)，
\begin{align*}
 \left(\frac{\omega_\lambda^2}{c^2}-\frac{\omega_{\lambda'}^2}{c^2}\right)
 \langle\mathbf f_\lambda,\mathbf f_{\lambda'}\rangle_\epsilon
 &=\langle\mathcal M\mathbf f_\lambda,\mathbf f_{\lambda'}\rangle_\epsilon
 -\langle\mathbf f_\lambda,\mathcal M\mathbf f_{\lambda'}\rangle_\epsilon\\
 &=\langle\mathbf f_\lambda,\mathcal M\mathbf f_{\lambda'}\rangle_\epsilon
 -\langle\mathbf f_\lambda,\mathcal M\mathbf f_{\lambda'}\rangle_\epsilon\\
 &=0,
\end{align*}
第二步用了自伴性 \(\langle\mathcal M\mathbf f_\lambda,\mathbf f_{\lambda'}\rangle_\epsilon=\langle\mathbf f_\lambda,\mathcal M\mathbf f_{\lambda'}\rangle_\epsilon\)。故 \(\langle\mathbf f_\lambda,\mathbf f_{\lambda'}\rangle_\epsilon=0\)，即模式在\eqnrefs{eq:maxwell-weighted-inner-product}下正交。这就是把 Maxwell 场化为独立正则模式所需的谱论基础。
\end{proof}

上述谱论对任意边界条件均成立，但最便于量子化的特例是真空周期盒——此时本征模取平面波，正交归一化退化为离散 $\delta$，场能直接分解为独立谐振子之和。

先把场限制在体积 \(V\) 的周期边界盒中。横向模式可以取为

\[
\mathbf u_{\mathbf k\lambda}(\mathbf r)
=\frac{\mathbf e_{\mathbf k\lambda}}{\sqrt V}
 \ee^{\ii\mathbf k\cdot\mathbf r},
\qquad
\mathbf k\cdot\mathbf e_{\mathbf k\lambda}=0,
\]

且

\[
\int_V\!\dd^3 r\,
\mathbf u_{\mathbf k\lambda}^*\cdot
\mathbf u_{\mathbf k'\lambda'}
=\delta_{\mathbf k\mathbf k'}\delta_{\lambda\lambda'}.
\]

把 \(\mathbf A\) 在这组模式上展开并代入经典场能
\[
H_{\rm EM}
=\frac12\int_V\!\dd^3 r\,
\left(\varepsilon_0\mathbf E^2+\frac{\mathbf B^2}{\mu_0}\right),
\]

可得到彼此独立的正则坐标 \(Q_{\mathbf k\lambda}\) 和共轭动量 \(P_{\mathbf k\lambda}\)：

\[
H_{\rm EM}
=\sum_{\mathbf k\lambda}
\left[
\frac{P_{\mathbf k\lambda}^2}{2}
+\frac{\omega_k^2Q_{\mathbf k\lambda}^2}{2}
\right],
\qquad \omega_k=c|\mathbf k|.
\]

因此``光场由谐振子组成''不是类比，而是 Maxwell 场在正交正常模上的正则分解。

\begin{proof}[解]
将矢势在正交模上展开为
\begin{equation*}
\mathbf A(\mathbf r,t)
=\sum_{\mathbf k\lambda}
Q_{\mathbf k\lambda}(t)\,\mathbf u_{\mathbf k\lambda}(\mathbf r),
\end{equation*}
其中 $Q_{\mathbf k\lambda}$ 是实正则坐标。电场和磁场由
\begin{equation*}
\mathbf E=-\frac{\partial\mathbf A}{\partial t}
=-\sum_{\mathbf k\lambda}\dot Q_{\mathbf k\lambda}\,\mathbf u_{\mathbf k\lambda},
\qquad
\mathbf B=\nabla\times\mathbf A
=\sum_{\mathbf k\lambda}Q_{\mathbf k\lambda}\,(\nabla\times\mathbf u_{\mathbf k\lambda})
\end{equation*}
给出。代入场能后，电场部分为
\begin{equation*}
\frac{\varepsilon_0}{2}\int_V\dd^3 r\,\mathbf E^2
=\frac{\varepsilon_0}{2}\sum_{\mathbf k\lambda}\sum_{\mathbf k'\lambda'}
\dot Q_{\mathbf k\lambda}\dot Q_{\mathbf k'\lambda'}
\int_V\dd^3 r\,\mathbf u_{\mathbf k\lambda}^*\cdot\mathbf u_{\mathbf k'\lambda'}
=\frac{\varepsilon_0}{2}\sum_{\mathbf k\lambda}\dot Q_{\mathbf k\lambda}^2,
\end{equation*}
其中最后一步用了正交关系 $\int_V\dd^3 r\,\mathbf u_{\mathbf k\lambda}^*\cdot\mathbf u_{\mathbf k'\lambda'}=\delta_{\mathbf k\mathbf k'}\delta_{\lambda\lambda'}$。磁场部分需利用横模满足的方程 $\nabla\times\nabla\times\mathbf u_{\mathbf k\lambda}=(\omega_k^2/c^2)\mathbf u_{\mathbf k\lambda}$，由此
\begin{equation*}
\int_V\dd^3 r\,(\nabla\times\mathbf u_{\mathbf k\lambda}^*)\cdot(\nabla\times\mathbf u_{\mathbf k'\lambda'})
=\int_V\dd^3 r\,\mathbf u_{\mathbf k\lambda}^*\cdot(\nabla\times\nabla\times\mathbf u_{\mathbf k'\lambda'})
=\frac{\omega_{k'}^2}{c^2}\delta_{\mathbf k\mathbf k'}\delta_{\lambda\lambda'},
\end{equation*}
其中第一步用了分部积分（周期边界使面积分消失），第二步用了模方程和正交性。因此磁场部分为
\begin{equation*}
\frac{1}{2\mu_0}\int_V\dd^3 rF\,\mathbf B^2
=\frac{1}{2\mu_0}\sum_{\mathbf k\lambda}
\frac{\omega_k^2}{c^2}Q_{\mathbf k\lambda}^2
=\frac{1}{2\mu_0\varepsilon_0 c^2}\sum_{\mathbf k\lambda}\omega_k^2 Q_{\mathbf k\lambda}^2
=\frac12\sum_{\mathbf k\lambda}\omega_k^2 Q_{\mathbf k\lambda}^2,
\end{equation*}
最后一步用了 $c^2=1/(\varepsilon_0\mu_0)$。定义共轭动量 $P_{\mathbf k\lambda}=\varepsilon_0\dot Q_{\mathbf k\lambda}$，电场部分化为 $\frac12\sum P_{\mathbf k\lambda}^2/\varepsilon_0$；若进一步取 $\varepsilon_0=1$ 的约定（或把 $P$ 重标为 $\sqrt{\varepsilon_0}\dot Q$），即得
\begin{equation*}
H_{\rm EM}=\sum_{\mathbf k\lambda}\left[\frac{P_{\mathbf k\lambda}^2}{2}+\frac{\omega_k^2 Q_{\mathbf k\lambda}^2}{2}\right],
\end{equation*}
即\eqnrefs{eq:maxwell-weighted-inner-product}所引用的谐振子分解。交叉项全部因正交性消失，每个 $(\mathbf k,\lambda)$ 模式独立贡献一个谐振子。
\end{proof}

经典场分解成独立正常模以后，对每个谐振子施加量子对易关系，就得到场算符；自发辐射率和一般环境中的 Green 并矢表示随后都从这个场算符推出。

\section{正则量子化、Green 张量与自发辐射}
施加正则对易关系
\[
[\hat Q_{\mathbf k\lambda},\hat P_{\mathbf k'\lambda'}]
=\ii\hbar\delta_{\mathbf k\mathbf k'}\delta_{\lambda\lambda'},
\]
并定义升降算符
\begin{equation}
\hat a_{\mathbf k\lambda}
=\sqrt{\frac{\omega_k}{2\hbar}}\left(\hat Q_{\mathbf k\lambda}
+\frac{\ii\hat P_{\mathbf k\lambda}}{\omega_k}\right),
\qquad
\hat a_{\mathbf k\lambda}^\dagger
=\sqrt{\frac{\omega_k}{2\hbar}}\left(\hat Q_{\mathbf k\lambda}
-\frac{\ii\hat P_{\mathbf k\lambda}}{\omega_k}\right),
\label{eq:ladder-operator-definition}
\end{equation}
其逆变换为 $\hat Q=\sqrt{\hbar/(2\omega)}\,((\hat a+\hat a^\dagger)$，$\hat P=-\ii\sqrt{\hbar\omega/2}\,(\hat a-\hat a^\dagger)$。由 $[\hat Q,\hat P]=\ii\hbar$ 直接验证 $[\hat a,\hat a^\dagger]=1$。将逆变换代入 $H=P^2/2+\omega^2 Q^2/2$：
\begin{align*}
\hat H
&=\frac12\left(-\frac{\hbar\omega}{2}\right)(\hat a-\hat a^\dagger)^2
+\frac{\omega^2}{2}\cdot\frac{\hbar}{2\omega}(\hat a+\hat a^\dagger)^2\\
&=\frac{\hbar\omega}{4}\left[-(\hat a^2-\hat a\hat a^\dagger-\hat a^\dagger\hat a+\hat a^{\dagger2})
+(\hat a^2+\hat a\hat a^\dagger+\hat a^\dagger\hat a+\hat a^{\dagger2})\right]\\
&=\frac{\hbar\omega}{4}\cdot 2(\hat a\hat a^\dagger+\hat a^\dagger\hat a)
=\frac{\hbar\omega}{2}(\hat a^\dagger\hat a+\hat a\hat a^\dagger)
=\hbar\omega\left(\hat a^\dagger\hat a+\frac12\right),
\end{align*}
其中最后一步用了 $[\hat a,\hat a^\dagger]=1$ 即 $\hat a\hat a^\dagger=\hat a^\dagger\hat a+1$。对每个模式独立执行此代换，得到
\[
\hat H_{\rm EM}
=\sum_{\mathbf k\lambda}
\hbar\omega_k
\left(
\hat a_{\mathbf k\lambda}^\dagger
\hat a_{\mathbf k\lambda}+\frac12
\right).
\]
其中 $\hat a_{\mathbf k\lambda}^\dagger\hat a_{\mathbf k\lambda}$ 是第 $(\mathbf k,\lambda)$ 模的光子数算符，$\frac12\hbar\omega_k$ 是真空零点能——对所有模式求和后形式上发散，但可观测的物理量只涉及差值（如 Lamb 移位、Casimir 力），发散部分在正常序 $\hat H=\sum\hbar\omega\,\hat a^\dagger\hat a$ 中自动消除。

将逆变换 $\hat Q_{\mathbf k\lambda}=\sqrt{\hbar/(2\omega_k)}\,(\hat a+\hat a^\dagger)$ 代入矢势展开 $\hat{\mathbf A}=\sum\hat Q_{\mathbf k\lambda}\mathbf u_{\mathbf k\lambda}$，并利用 $\mathbf u_{\mathbf k\lambda}=\mathbf e_{\mathbf k\lambda}\ee^{\ii\mathbf k\cdot\mathbf r}/\sqrt V$，得到
\[
\hat{\mathbf A}(\mathbf r,t)
=\sum_{\mathbf k\lambda}
\sqrt{\frac{\hbar}{2\varepsilon_0\omega_kV}}
\left[
\mathbf e_{\mathbf k\lambda}
\hat a_{\mathbf k\lambda}
 \ee^{\ii(\mathbf k\cdot\mathbf r-\omega_kt)}
+\text{H.c.}
\right].
\]
其中前因子 $\sqrt{\hbar/(2\varepsilon_0\omega_kV)}$ 来自 $\sqrt{\hbar/(2\omega_k)}\cdot 1/\sqrt V$ 再乘 $\varepsilon_0^{-1/2}$（由 $P=\varepsilon_0\dot Q$ 约定引入的 $\varepsilon_0$ 归一化）。

于是正频电场为

\begin{equation}
\hat{\mathbf E}^{(+)}(\mathbf r,t)
=\ii\sum_{\mathbf k\lambda}
\sqrt{\frac{\hbar\omega_k}{2\varepsilon_0V}}
\mathbf e_{\mathbf k\lambda}
\hat a_{\mathbf k\lambda}
 \ee^{\ii(\mathbf k\cdot\mathbf r-\omega_kt)},
\label{eq:field-expansion}
\end{equation}

而 \(\hat{\mathbf E}^{(-)}=[\hat{\mathbf E}^{(+)}]^\dagger\)。连续自由空间极限通过 \(V^{-1}\sum_{\mathbf k}\to(2\pi)^{-3}\int\dd^3 k\) 得到。

场算符一旦建立，自发辐射便不再是唯象引入的 Einstein $A$ 系数，而是原子偶极与真空涨落耦合的直接结果。以下从 Fermi 黄金法则出发，将 $A_{eg}$ 完全用偶极矩阵元和频率表达。

Einstein \(A\) 系数在\chapref{ch:time-dependent-perturbation}中是通过热平衡关系引入的。量子化电磁场使之可直接写成偶极矩阵元。考虑初态

\[
|i\rangle=|e;0\rangle
\]

和末态

\[
|f\rangle=|g;1_{\mathbf k\lambda}\rangle.
\]

在电偶极近似中

\[
H_I=-\hat{\mathbf d}\cdot\hat{\mathbf E}.
\]

只有负频部分 \(\hat{\mathbf E}^{(-)}\propto a^\dagger\) 能从真空产生一个光子，因此

\[
M_{\mathbf k\lambda}
\equiv
\langle g;1_{\mathbf k\lambda}|H_I|e;0\rangle
=-\ii
\sqrt{\frac{\hbar\omega_k}{2\varepsilon_0V}}
\,\mathbf d_{ge}\cdot\mathbf e_{\mathbf k\lambda}^*
 \ee^{-\ii\mathbf k\cdot\mathbf r_0}.
\]

Fermi 黄金法则给出

\[
\Gamma_{e\to g}
=\frac{2\pi}{\hbar}
\sum_{\mathbf k\lambda}
|M_{\mathbf k\lambda}|^2
\delta(E_g+\hbar\omega_k-E_e).
\]

\begin{proof}[解]
矩阵元。由 $\hat{\mathbf E}^{(-)}=-\ii\sum_{\mathbf k'\lambda'}\sqrt{\hbar\omega_{k'}/(2\varepsilon_0 V)}\,\mathbf e_{\mathbf k'\lambda'}^*\hat a_{\mathbf k'\lambda'}^\dagger\ee^{-\ii(\mathbf k'\cdot\mathbf r-\omega_{k'}t)}$，矩阵元为
\begin{equation*}
M_{\mathbf k\lambda}
=-\langle g;1_{\mathbf k\lambda}|\hat{\mathbf d}\cdot\hat{\mathbf E}^{(-)}|e;0\rangle
=\ii\sum_{\mathbf k'\lambda'}\sqrt{\frac{\hbar\omega_{k'}}{2\varepsilon_0 V}}\,\mathbf d_{ge}\cdot\mathbf e_{\mathbf k'\lambda'}^*\ee^{-\ii\mathbf k'\cdot\mathbf r_0}\,\langle 1_{\mathbf k\lambda}|\hat a_{\mathbf k'\lambda'}^\dagger|0\rangle,
\end{equation*}
其中 $\mathbf d_{ge}=\langle g|\hat{\mathbf d}|e\rangle$ 是原子偶极矩阵元，时间因子在共振频率 $\omega_0=(E_e-E_g)/\hbar$ 处取值。由 $\hat a_{\mathbf k'\lambda'}^\dagger|0\rangle=|1_{\mathbf k'\lambda'}\rangle$ 和正交性 $\langle 1_{\mathbf k\lambda}|1_{\mathbf k'\lambda'}\rangle=\delta_{\mathbf k\mathbf k'}\delta_{\lambda\lambda'}$，求和坍缩为单一项，即 $M_{\mathbf k\lambda}$ 的表达式。

模平方。取 $|M_{\mathbf k\lambda}|^2$，相位因子 $\ee^{-\ii\mathbf k\cdot\mathbf r_0}$ 的模为 1，因此
\begin{equation*}
|M_{\mathbf k\lambda}|^2
=\frac{\hbar\omega_k}{2\varepsilon_0 V}
|\mathbf d_{ge}\cdot\mathbf e_{\mathbf k\lambda}|^2.
\end{equation*}

连续极限与偏振求和。把 $\sum_{\mathbf k}\to V/(2\pi)^3\int\dd^3 k$ 代入 Fermi 黄金法则，$V$ 消去：
\begin{equation*}
\Gamma=\frac{2\pi}{\hbar}\sum_\lambda\frac{V}{(2\pi)^3}\int\dd^3 k\,\frac{\hbar\omega_k}{2\varepsilon_0 V}|\mathbf d\cdot\mathbf e_{\mathbf k\lambda}|^2\delta(E_g+\hbar\omega_k-E_e).
\end{equation*}
利用 $\delta(E_g+\hbar\omega-E_e)=\delta(\omega-\omega_0)/\hbar$，对 $\mathbf k$ 用球坐标 $\dd^3 k=k^2\dd k\dd\Omega$，先做偏振求和
\begin{equation*}
\sum_{\lambda=1,2}|\mathbf d\cdot\mathbf e_{\mathbf k\lambda}|^2
=|\mathbf d|^2-|\mathbf d\cdot\hat{\mathbf k}|^2,
\end{equation*}
这是横向偏振完备关系：两个偏振矢量 $\mathbf e_{\mathbf k1},\mathbf e_{\mathbf k2}$ 与 $\hat{\mathbf k}$ 构成三维正交基，$\sum_\lambda|\mathbf d\cdot\mathbf e_{\mathbf k\lambda}|^2=\mathbf d^\dagger(\mathbf I-\hat{\mathbf k}\hat{\mathbf k}^T)\mathbf d$。

角积分。取 $\mathbf d$ 沿 $z$ 轴，则 $|\mathbf d\cdot\hat{\mathbf k}|^2=|\mathbf d|^2\cos^2\theta$，
\begin{equation*}
\int\dd\Omega\,(|\mathbf d|^2-|\mathbf d|^2\cos^2\theta)
=|\mathbf d|^2\cdot 2\pi\int_0^\pi\sin\theta\,\sin^2\theta\,\dd\theta
=|\mathbf d|^2\cdot 2\pi\cdot\frac{4}{3}
=\frac{8\pi}{3}|\mathbf d|^2.
\end{equation*}

径向积分。由 $k=\omega/c$ 得 $k^2\dd k=\omega^2\dd\omega/c^3$，\(\delta\) 函数 $\delta(\omega-\omega_0)$ 将积分锁定在 $\omega_0$：
\begin{equation}
\label{eq:spontaneous-emission-rate}
\Gamma=\frac{2\pi}{\hbar}\cdot\frac{1}{(2\pi)^3}\cdot\frac{\hbar}{2\varepsilon_0}\cdot\frac{8\pi}{3}|\mathbf d|^2\cdot\frac{\omega_0^2}{c^3}\cdot\omega_0\cdot\frac{1}{\hbar}
=\frac{\omega_0^3|\mathbf d|^2}{3\pi\varepsilon_0\hbar c^3},
\end{equation}
即\eqnrefs{eq:spontaneous-emission-rate}。其中 $\omega_0^3$ 的物理来源：一个 $\omega_0$ 来自单光子电场振幅 $\sqrt{\omega}$ 的模平方，两个 $\omega_0^2$ 来自三维态密度 $k^2\dd k=\omega^2\dd\omega/c^3$。
\end{proof}

这说明自发辐射不是经典外场触发的跃迁，而是原子偶极与真空电磁场连续模耦合的结果。\(\omega_0^3\) 来自单光子场振幅中的 \(\sqrt\omega\) 与三维光子密度 of 态中的 \(\omega^2\) 共同作用。把 \eqnrefs{eq:spontaneous-emission-rate} 与 Planck 谱结合，即可恢复 Einstein \(A\)-\(B\) 关系。

\eqnrefs{eq:spontaneous-emission-rate} 对氢原子 \(2p\to1s\) 跃迁（Lyman-\(\alpha\)，\(\lambda=121.6\,\mathrm{nm}\)，\(\omega_0=1.55\times10^{16}\,\mathrm{rad/s}\)，\(|\mathbf d_{ge}|=ea_0/\sqrt{2}\approx6.0\times10^{-30}\,\mathrm{C\,m}\)）：
\[
A_{2p\to1s}
=\frac{\omega_0^3|\mathbf d|^2}{3\pi\epsilon_0\hbar c^3}
\approx6.27\times10^8\,\mathrm{s^{-1}}
\]
对应寿命 \(\tau=1/A\approx1.6\,\mathrm{ns}\)，与实验值 \(1.60\,\mathrm{ns}\) 精确吻合。对比钠原子 \(3p\to3s\) 跃迁（D 线，\(\lambda=589\,\mathrm{nm}\)，\(|\mathbf d|\approx2.5\times10^{-29}\,\mathrm{C\,m}\)）：\(A_{3p\to3s}\approx6.1\times10^7\,\mathrm{s^{-1}}\)，\(\tau\approx16\,\mathrm{ns}\)，与实验值 \(16.3\,\mathrm{ns}\) 吻合。\(\omega^3\) 依赖使紫外跃迁的辐射率远大于可见光跃迁。

自由空间的结果\eqnrefs{eq:spontaneous-emission-rate}依赖逐模求和，但在腔、界面或光子晶体等非均匀环境中，直接求模并不方便。Green 并矢提供了一个统一框架：其虚部编码局域光子态密度，自发辐射率只需在 $\mathbf r_0$ 处求值即可。

在非均匀无耗散环境中，逐个求模并非总是最方便。定义出射 electromagnetic Green 并矢
\begin{equation}
\left[\nabla\times\nabla\times
-\frac{\omega^2}{c^2}\epsilon_r(\mathbf r)\right]
\mathbf G(\mathbf r,\mathbf r';\omega)
=\mathbf I\,\delta^{(3)}(\mathbf r-\mathbf r').
\label{eq:maxwell-green-dyadic}
\end{equation}
其虚部编码局域光子态密度。位于 $\mathbf r_0$ 的电偶极跃迁在弱耦合、马尔可夫极限下具有一般自发辐射率
\begin{equation}
\Gamma(\mathbf r_0)
=\frac{2\omega_0^2}{\hbar\epsilon_0c^2}
\mathbf d_{ge}^*\cdot
\operatorname{Im}\mathbf G(\mathbf r_0,\mathbf r_0;\omega_0)
\cdot\mathbf d_{ge}.
\label{eq:spontaneous-green-dyadic}
\end{equation}
自由空间有
\[
\operatorname{Im}\mathbf G_0(\mathbf r_0,\mathbf r_0;\omega)
=\frac{\omega}{6\pi c}\mathbf I,
\]
故\eqnrefs{eq:spontaneous-green-dyadic}立即退化为\eqnrefs{eq:spontaneous-emission-rate}。腔 Purcell enhancement、界面附近抑制/增强和光子晶体带隙效应，本质上都是 $\operatorname{Im}\mathbf G$ 改变，而不是“原子的 $A$ 系数本身神秘变化”。

\begin{proof}[解]
对\eqnrefs{eq:maxwell-normal-mode-eigenproblem}的正交本征模 $\{\mathbf f_\lambda\}$，满足 $\int\dd^3r\,\epsilon_r(\mathbf r)\,\mathbf f_\lambda\cdot\mathbf f_{\lambda'}^*=\delta_{\lambda\lambda'}$，Green 并矢的谱表示为
\begin{equation*}
\mathbf G(\mathbf r,\mathbf r';\omega)
= c^2\sum_\lambda
 \frac{\mathbf f_\lambda(\mathbf r)\,\mathbf f_\lambda^*(\mathbf r')}
 {\omega_\lambda^2-(\omega+\ii0^+)^2}.
\end{equation*}
取 $\mathbf r=\mathbf r'=\mathbf r_0$ 并对分母因式分解，
\begin{equation*}
\omega_\lambda^2-(\omega+\ii0^+)^2
= (\omega_\lambda-\omega-\ii0^+)(\omega_\lambda+\omega+\ii0^+).
\end{equation*}
在正频 $\omega_0>0$ 附近，第二因子近似 $2\omega_0$，因此
\begin{equation*}
\mathbf G(\mathbf r_0,\mathbf r_0;\omega_0)
\approx \frac{c^2}{2\omega_0}\sum_\lambda
\frac{\mathbf f_\lambda(\mathbf r_0)\,\mathbf f_\lambda^*(\mathbf r_0)}
{\omega_\lambda-\omega_0-\ii0^+}.
\end{equation*}
利用 Sokhotski--Plemelj 公式
\begin{equation*}
\operatorname{Im}\frac{1}{x-\ii0^+}=\pi\,\delta(x),
\end{equation*}
取虚部，
\begin{equation*}
\operatorname{Im}\mathbf G(\mathbf r_0,\mathbf r_0;\omega_0)
= \frac{\pi c^2}{2\omega_0}\sum_\lambda
 \mathbf f_\lambda(\mathbf r_0)\,\mathbf f_\lambda^*(\mathbf r_0)\,
 \delta(\omega_0-\omega_\lambda).
\end{equation*}
另一方面，量子化后单光子电场矩阵元为
\begin{equation*}
\langle 0|\hat{\mathbf E}(\mathbf r_0)|1_\lambda\rangle
= \ii\sqrt{\frac{\hbar\omega_\lambda}{2\epsilon_0}}\,
 \mathbf f_\lambda(\mathbf r_0),
\end{equation*}
代入 Fermi 黄金法则
\begin{equation*}
\Gamma = \frac{2\pi}{\hbar}\sum_\lambda
 |\mathbf d_{ge}\cdot\langle 0|\hat{\mathbf E}(\mathbf r_0)|1_\lambda\rangle|^2
 \delta(\hbar\omega_0-\hbar\omega_\lambda),
\end{equation*}
展开模平方并把 $\delta(\hbar\omega_0-\hbar\omega_\lambda)=\delta(\omega_0-\omega_\lambda)/\hbar$ 代入，模式求和的结构与 $\operatorname{Im}\mathbf G$ 完全一致，因此
\begin{equation*}
\Gamma(\mathbf r_0)
= \frac{2\omega_0^2}{\hbar\epsilon_0 c^2}\,
 \mathbf d_{ge}^*\cdot
 \operatorname{Im}\mathbf G(\mathbf r_0,\mathbf r_0;\omega_0)\cdot
 \mathbf d_{ge},
\end{equation*}
即\eqnrefs{eq:spontaneous-green-dyadic}。Green 并矢是"模式密度乘局域场强"的连续谱压缩表示。
\end{proof}
